\section{Literature Review Procedure}
\label{sec:lit-review}
The literature review has two tasks in this study. It collects the harms and capabilities used
to construct harm pathways, and it supplies the evidence for the methodological decisions of
this chapter. Its procedure follows the staged review of Baumann et
al.~\parencite{baumann2025llm}: fixed Boolean queries built from term groups, the exact string
and hit count recorded for every source (Appendix~\ref{app:lit-strings}), removal of
duplicates, screening against explicit criteria and the extraction of fixed fields. The search is reported with the items of the
PRISMA-S extension and the flow of records with the flow diagram of
PRISMA~2020~\parencite{rethlefsen2021prisma,page2021prisma}. Because the review maps evidence
rather than pooling effects, it is a scoping review in the sense of Arksey and
O'Malley~\parencite{arksey2005scoping}. The queries and criteria are fixed before the search.
\highlight[orange!25]{\textbf{[?]}}\todo{Missing (PRISMA-ScR item 5, PRISMA-S item 14): say
whether a review protocol is registered and whether a second person reviews the search
strategies, e.g. \enquote{No protocol is registered, and no second searcher reviews the
strategies.}}

\subsection{Review Questions and Search Concepts}
\label{sec:lit-search}
\paragraph{Review questions.} The review has four sub-questions, derived from the research
question \enquote{which harms health-tracker providers communicate}, each tied to a part of
this study.
\begin{description}[leftmargin=1.5em,itemsep=1pt,topsep=3pt,labelindent=0pt,font=\normalfont\bfseries]
  \item[LQ1] \textit{Harms and capabilities.} Which privacy harms, risks, threats
    and inference capabilities concern the users of health tracking apps and
    their wearables? It supplies evidence for constructing the pathways assessed in RQ1
    and RQ3, and starts from general accounts of privacy harm.
  \item[LQ2] \textit{Communication.} How has the privacy communication of health
    tracking apps been analysed, and what was found? It informs
    the sampling, material and coding decisions and tests the gap for RQ1--RQ4.
  \item[LQ3] \textit{Language models.} How have large language models been used
    to annotate, classify, extract from or assess privacy text, and how were
    they validated? It informs the LLM-assisted coding.
  \item[LQ4] \textit{Framing.} How do privacy texts in any domain word or present
    risks, harms and protections, for example through vagueness, hedging,
    reassurance or misleading claims? It informs the stance and protection axes
    (RQ2, RQ3).
\end{description}

\paragraph{Search concepts.} The search is built from concept blocks. Each
concept of the sub-questions becomes a term group whose alternative terms are
joined with OR, and the groups of a query are joined with AND, so that a record
must contain at least one term of every group. Eight groups cover the concepts
(Table~\ref{tab:lit-groups}; the complete term lists are in
Appendix~\ref{app:lit-groups}).

\begin{table}[H]
\centering
\begin{singlespace}
\caption[Term groups of the literature search]{Term groups of the literature
search, with the number of terms and examples; the complete lists are in
Table~\ref{tab:lit-groups-full}. The terms of a group are joined with OR. A
trailing * truncates, and terms of several words are searched as phrases.
Table~\ref{tab:lit-strands} shows how the groups are combined.}
\label{tab:lit-groups}
\small
\begin{tabularx}{\textwidth}{@{}l>{\raggedright\arraybackslash}p{3.6cm}r>{\raggedright\arraybackslash}X@{}}
\toprule
\textbf{Group} & \textbf{Concept} & \textbf{Terms} & \textbf{Examples} \\
\midrule
\input{tables/literature-groups-short}% the generated rows end with \bottomrule
\end{tabularx}
\end{singlespace}
\end{table}

Four strands combine the groups, one for each sub-question
(Table~\ref{tab:lit-strands}). The strands cover 1~January 2014 to the
search date, which includes the first studies of health-app privacy
policies~\parencite{sunyaev2015availability}; S3 starts in 2022, as in the
review of language-model studies by Baumann et al.

\begin{table}[H]
\centering
\begin{singlespace}
\caption[Search strands]{Search strands, one for each sub-question. A dot marks
a group that the strand requires; the marked groups are joined with AND.}
\label{tab:lit-strands}
\small
\begin{tabular*}{\textwidth}{@{\extracolsep{\fill}}ll*{8}{c}lr@{}}
\toprule
 & & \multicolumn{8}{c}{\textbf{Term group}} & & \\
\cmidrule(lr){3-10}
\textbf{Strand} & \textbf{Sub-question} & D & P & H & C & T & L & A & F & \textbf{Query} & \textbf{Time frame} \\
\midrule
S1 & LQ1 Harms and capabilities & $\bullet$ & $\bullet$ & $\bullet$ & & & & & & D AND P AND H & 2014--2026 \\
S2 & LQ2 Communication & $\bullet$ & $\bullet$ & & $\bullet$ & & & & & D AND P AND C & 2014--2026 \\
S3 & LQ3 Language models & & & & & $\bullet$ & $\bullet$ & $\bullet$ & & T AND L AND A & 2022--2026 \\
S4 & LQ4 Framing & & & & & $\bullet$ & & & $\bullet$ & T AND F & 2014--2026 \\
\bottomrule
\end{tabular*}
\end{singlespace}
\end{table}

\subsection{Sources, Search Execution, and Citation Searching}
\label{sec:lit-sources}
\paragraph{Venues.} The eight venues of security, privacy and human--computer
interaction research are each listed completely from 2014 to the search date, apart from the
gaps that Table~\ref{tab:lit-venues} gives. They are the four main security conferences (CCS, IEEE
S\&P, NDSS and USENIX Security), the venues for privacy-enhancing technologies
(PoPETs) and for usable privacy and security (SOUPS), and CHI and CSCW, where
user studies of privacy in health tracking appear. The last column of the table
gives the selection rule of each venue. At the security and privacy venues, privacy is
the subject of the venue itself, so the domain group alone selects the
candidates for LQ1 and LQ2.

\paragraph{Databases.} Four bibliographic databases add the literature published
elsewhere (Table~\ref{tab:lit-databases}): PubMed for health informatics and
medicine, IEEE Xplore for engineering, arXiv for preprints and the ACL Anthology
for language processing. Each strand is run against every database, as a query
of its own or locally on the retrieved records.

\paragraph{Citation searching.} Citation searching adds studies that the term
groups miss. It starts from a seed set of four studies that address the core
topic of the review directly: three CHI papers on privacy in period, fertility
and other female health
apps~\parencite{malki2024exploring,hudig2025intimate,cao2024deleted} and a
survey of the utility, privacy and security of wearable activity
trackers~\parencite{niksirat2024wearable}. One round of backward searching (the
references of the seed papers) and forward searching (the works that cite them)
is run on \highlight[orange!25]{\textbf{[?]}}\todo{Missing (PRISMA-S item 13): the date of
the citation search.}. References and citing works are taken from the OpenCitations Index and
their metadata from Crossref. The records found enter the same deduplication
and screening as the search results.

\subsection{Screening and Deduplication}
\label{sec:lit-screening}
\paragraph{Deduplication.} The records of all sources are merged automatically
by a script of this project. Two records are taken to be the same work when
they share a DOI or an arXiv identifier, or when their normalised titles are identical and
their publication years differ by at most one (three if one of them is an arXiv preprint).
The merged record keeps every source that found it.

\paragraph{Eligibility.} A record is eligible if it meets all four criteria of
Table~\ref{tab:lit-eligibility}, which were fixed before screening.

\begin{table}[H]
\centering
\begin{singlespace}
\caption[Eligibility criteria of the literature review]{Eligibility criteria of
the literature review.}
\label{tab:lit-eligibility}
\small
\begin{tabularx}{\textwidth}{@{}l>{\raggedright\arraybackslash}p{2.6cm}>{\raggedright\arraybackslash}X@{}}
\toprule
 & \textbf{Criterion} & \textbf{Eligibility requirement} \\
\midrule
C1 & Publication type & A journal article, conference paper, workshop paper, or archival preprint that reports original empirical or analytical research or presents a systematic or scoping review. \\
C2 & Language & Written in English. \\
C3 & Publication date & Published between 1 January 2014 and 28 October 2026. For LQ3, the publication window begins in 2022. \\
C4 & Relevance & Addresses at least one of the four literature review questions (LQ1--LQ4). \\
\bottomrule
\end{tabularx}
\end{singlespace}
\end{table}

\paragraph{Screening.} The exclusion rules of Table~\ref{tab:lit-criteria} were
fixed before screening and are applied in their order, at the stage the table
gives; the abstract is read where the title does not decide. Every decision and
its reason are recorded in a screening log. The decisions are made by
\highlight[orange!25]{\textbf{[?]}}.\todo{Missing (PRISMA 2020 item 8, PRISMA-ScR item 9): who
screens, how many people, and whether an AI assistant (name, version) prepares the decisions.
The old text read: \enquote{The screening decisions were prepared by an AI assistant from
titles, abstracts and full texts. The author checks every inclusion and every full-text
exclusion, and a random fifth of the title-and-abstract exclusions drawn with a recorded seed;
if that fifth holds a single eligible record, every title-and-abstract exclusion is checked.}
Use it only once that check is done.}
Where no source holds an abstract, the abstract of the same DOI is taken from
OpenAlex~\parencite{priem2022openalex}. Two conventions keep the rules apart:
systems for hospitals, health records, telemedicine or the Internet of Things
in general fall under R3, and systems for wearables, body-area networks or
health apps under R5; adoption, attitude and concern studies fall under R4
unless they name specific harms or actors. The review maps evidence, so no quality
score decides inclusion, as PRISMA-ScR allows for scoping
reviews~\parencite{tricco2018prisma}; the peer-review status of each study is
recorded instead, and preprints are marked as such.

\begin{table}[H]
\centering
\begin{singlespace}
\caption[Exclusion rules of the literature review]{Exclusion rules of the
literature review. A record is excluded under the first rule that applies.}
\label{tab:lit-criteria}
\small
\begin{tabularx}{\textwidth}{@{}l>{\raggedright\arraybackslash}X>{\raggedright\arraybackslash}p{2.0cm}@{}}
\toprule
 & \textbf{Rule} & \textbf{Stage} \\
\midrule
R1 & Not a research article: editorial, commentary, letter, erratum, abstract, poster, keynote, book or chapter, thesis, front matter, narrative overview, position paper & title, abstract \\
R2 & Outside the language or period of C2--C3 & title, abstract \\
R3 & Outside the domain: neither health tracking apps or wearables nor privacy text (for example hospital systems, health records, telemedicine, clinical sensor networks) & title, abstract \\
R4 & Health tracking without a privacy focus (quality, usability, clinical effect, adoption, features) & title, abstract \\
R5 & A technical system or protocol whose evaluation documents neither harms or capabilities nor communication & title, abstract \\
R6 & Privacy text or language models outside health tracking, answering neither LQ3 nor LQ4 (for example models as a privacy risk, readability, compliance or label accuracy without risk wording) & title, abstract \\
R7 & Full text requires payment and is available neither through the university's institutional access or network nor as a publicly accessible version (for example an arXiv preprint) & full text \\
R8 & On the full text, C1 or C4 does not hold & full text \\
\bottomrule
\end{tabularx}
\end{singlespace}
\end{table}

\subsection{Extraction and Synthesis}
\label{sec:lit-synthesis}
Each study that passes title-and-abstract screening receives one row in an extraction table.
The row records the review sub-question assigned during screening and three fields coded from
the title, abstract and venue using fixed keyword lists: app domain, material studied and venue
type.

For included studies, manual full-text coding by
\highlight[orange!25]{\textbf{[?]}}\todo{Missing (PRISMA-ScR item 10): who charts the full
texts, and whether a second person verifies the extraction.} records which of RQ1--RQ4 the study informs or
whether it provides a methodological precedent, the harms it names and, for LQ1 studies, the
capabilities it documents. Capabilities are grouped into eight categories (K1--K8), each
recorded with the data it requires and the actor who could use it. The resulting capability
table informs the construction of harm pathways, corresponding to Baumann et al.'s third stage.

The descriptive synthesis counts studies by year and venue, maps evidence and gaps across app
domains, materials and research questions, and summarises findings for each review
sub-question.

